\par\Needspace{4\baselineskip}
\begin{center}
\hypertarget{native-036-article018}{}\pdfbookmark[2]{18.}{article-018}18.
\end{center}
In the PARABOLA, the eccentric anomaly, the mean anomaly and the mean motion would all be $=0$. These notions therefore cannot serve here to relate the motion to time. In the parabola, however, we have no need of an auxiliary angle to integrate $r r\,\mathrm{d}v$, for
\[
\begin{aligned}
r r\,\mathrm{d}v&=\frac{p p\,\mathrm{d}v}{4\cos\frac{1}{2}v^4}=\frac{p p\,\mathrm{d}\operatorname{tang}\frac{1}{2}v}{2\cos\frac{1}{2}v^2}\\
&=\frac{1}{2}p p(1+\operatorname{tang}\frac{1}{2}v^2)\,\mathrm{d}\operatorname{tang}\frac{1}{2}v
\end{aligned}
\]
and hence $\int r r\,\mathrm{d}v=\frac{1}{2}p p(\operatorname{tang}\frac{1}{2}v+\frac{1}{3}\operatorname{tang}\frac{1}{2}v^3)+\mathrm{Const.}$ If time is reckoned from the passage through perihelion, the constant is $=0$. We therefore have
\[
\operatorname{tang}\frac{1}{2}v+\frac{1}{3}\operatorname{tang}\frac{1}{2}v^3=\frac{2t k\BareRoot{(1+\mu)}}{p^{\frac{3}{2}}},
\]
\SourceLeaf{037}which gives $t$ in terms of $v$, and $v$ in terms of $t$, as soon as $p$ and $\mu$ are known. In place of $p$, it is customary to give among the parabolic elements the radius vector at perihelion, which is $\frac{1}{2}p$, and to neglect the mass $\mu$ altogether. It is very unlikely that we shall ever be able to determine the mass of a body whose orbit is computed as a parabola. Indeed, according to the best and most recent observations, all comets appear to have so small a density and mass that the mass may be regarded as imperceptible and safely neglected.
